\section{Methodology}

% Source: .context/chat/Paper Overview (2026-11-19) (2).md.
% Structure: nine paragraph roles plus two foregrounded formal results;
% provisional drafting budget 1,400--1,800 words.
% Governing distinction: information recovery, financial geometry and
% spatial-economic consequences are separate claims.

% P1 -- Economic and statistical objects.
% Define firms, dates, returns, factors, accounting histories and dated
% information sets; distinguish issuer and security units.
% Define target explicitly: factor exposures, residual dependence or another
% economic quantity; never switch targets between theorem and experiment.
% Separate Fernandez characteristics from expanded accounting inputs.

% P2 -- Conventional characteristic distances.
% Preserve Fernandez-style historical-ratio correlation distances and weights.
% Prespecify standardized Euclidean, Mahalanobis and PCA comparators if the
% active study supports them.
% State overlap, undefined-correlation, bandwidth and normalization rules.
% US reimplementation kernel behavior cannot explain unavailable original
% sample estimates.

% P3 -- Frozen TFM representation and estimation inventory.
% Define fixed encoder E_{theta_0}, row embeddings and model-specific context.
% Inventory preprocessing, reference context, exposure readout, factor
% covariance, innovation variances, spatial coefficient and distance
% calibration.
% Distinguish native fixed features from financially supervised calibration.
% Freezing removes encoder weight estimation on the financial sample, but does
% not confer automatic admissibility.
% SPOTTING -- Dependence/complexity diagram or controlled learning curve.

% P4 -- Native and calibrated geometry.
% Define cosine, angular, chordal and decoded-exposure distances.
% Chordal construction: primary PCoA-compatible candidate where applicable.
% Compare native, exposure-calibrated and metric-calibrated procedures with
% controls for normalization, dimension and supervision where possible.
% Distinguish recoverable information from native nearest neighbors.

% P5 -- Theorem 1: representability without geometric identification.
% State target recovery through constancy on encoder fibers.
% Give an invertible non-isometric transformation that changes distances while
% preserving recovery; define the transformation class precisely.
% Separate exact from approximate recovery and sufficient certificates from
% necessary conditions.
% Interpretation: limit on native-distance claims, not proof all
% embeddings fail.
% Put measurability detail and full Lean mapping in appendices.

% P6 -- Spatial weight construction.
% Define affinity kernel, bandwidth, diagonal treatment, sparsification and
% normalization for the actual implemented network.
% Report symmetry, spectrum, row sums and degree ratios where applicable.
% Similarity networks are not directed causal transmission networks.
% State identification and multiplier-stability conditions for each model.

% P7 -- Spatial CAPM and S-APT specifications.
% Present factor benchmark, SAR, SLX and SEM with model-specific reduced-form
% exposures and residual covariance.
% State inherited pricing restrictions and maintained factor, trading and
% no-asymptotic-arbitrage assumptions separately from covariance fit.
% Distinguish pricing errors, covariance fit and spatial-coefficient inference.
% Explain what the active APT residual-covariance experiment identifies.

% P8 -- Theorem 2: spectral and financial-risk stability.
% Use the implemented normalization; state a sufficient spectral gap and
% bounded-innovation premise for multiplier and covariance control.
% Trace distance/embedding error -> network error -> multiplier error ->
% covariance error -> portfolio-risk error with explicit scope.
% Bounded multiplier or positive-definite covariance is not a
% no-arbitrage proof.
% Place finite-probability constructions and intermediate lemmas in appendices.

% P9 -- Inference and economic interpretation.
% Specify uncertainty for generated weights, factors, spatial parameters and
% pricing contrasts; distinguish conditional from full uncertainty.
% Use dependence-aware intervals, fit diagnostics and multiplicity control.
% Define PCoA alignment, subspace stability, factor-mimicking construction and
% attribution as distinct interpretation tests.
% Risk endpoints: in-sample covariance and model-implied VaR under stated
% distribution assumptions; historical fit, not forecasting claims.
